% §9.5 Bridge to Sketches
%
% Monograph register: compact forward transition to chapter 10. Cross-refs
% def:n-ary-channel, ax:element-indep, ax:block-rates-indep, prop:n-ary-capacity,
% ch:sketches. DEFINED: sec:bridge-to-sketches.

\section{Bridge to Sketches}
\label{sec:bridge-to-sketches}

A \emph{sketch} maps a stream or multiset to a compact summary, often sublinear
in the input size, from which a query routine reconstructs an answer with error
bounded by an analytic expression in the summary size: the user trades a known
accuracy for a known space budget. What distinguishes a sketch from lossy
compression is that bounded, calculable error. The three canonical sketches,
CountMin for frequency, HyperLogLog for cardinality, and MinHash for similarity,
are each a $\Bernoulli{T}$ instance of this chapter.

Each is a channel in the sense of \Cref{def:n-ary-channel}, differing only in
the structure of $Q$. CountMin is a counter ($T = \mathbb{N}$) whose one-sided
overestimation puts all of $Q$'s mass at or above the diagonal; HyperLogLog is a
counter whose mass concentrates in a multiplicative band around the diagonal
(relative rather than one-sided error); MinHash is a similarity estimator
($T = [0,1]$) whose $Q$ is a continuous kernel with output concentrated in a
band of width $O(1/\sqrt{r})$ at sketch size $r$. All three share the same
skeleton, an input type, an output type, a row-stochastic kernel, and the
per-query independence of Axioms~1 and~2 (\Cref{ax:element-indep},
\Cref{ax:block-rates-indep}); each one's error bound is a statement about the
entries of $Q$, its composition rule a statement about how $Q$-matrices combine,
and its capacity a specialization of \Cref{prop:n-ary-capacity}.
\Cref{ch:sketches-estimators} builds the three from scratch, the first non-Boolean
instances of $\Bernoulli{T}$ the book treats in detail.
